\Needspace{14\baselineskip}
\SeventhArticle{18}
The investigations of the preceding article do show how to find suitable functions $U$ for each value of the number $n$, but only successively, since one must pass from smaller values to larger ones. We readily observe, however, that these functions are expressed in general by
\[
\begin{aligned}&
u^{n+1}-\frac{(n+1)n}{2\mathbin{.}(2n+1)}u^{n-1}
\\[5pt]&+\frac{(n+1)n(n-1)(n-2)}{2\mathbin{.}4(2n+1)(2n-1)}u^{n-3}
\\[5pt]&-\frac{(n+1)n(n-1)(n-2)(n-3)(n-4)}{2\mathbin{.}4\mathbin{.}6\mathbin{.}(2n+1)(2n-1)(2n-3)}u^{n-5}\\ &
+\text{etc.}
\end{aligned}
\]
or also, if we use the symbol $F$ as in the paper cited above, by
\[
u^{n+1}F(-\tfrac12n,\ -\tfrac12(n+1),\ -(n+\tfrac12),\ u^{-2})
\]
This induction can easily be turned into a rigorous proof by the well-known method, or, if one prefers, with the aid of formula 19 in the cited paper. The function $U$, if one prefers, can also be expressed with the order of the terms reversed, by
\[
\pm\frac{3\mathbin{.}5\mathbin{.}7\mathinner{.\,.\,.\,.}(n+1)}{(n+3)(n+5)\mathinner{.\,.\,.\,.}(2n+1)}\mathbin{.}uF(-\tfrac12n,\ \tfrac12(n+3),\ \tfrac32,\ uu)
\]
for even $n$, with the upper or lower sign according as $\frac12n$ is even or odd; or by
\[
\pm\frac{1\mathbin{.}3\mathbin{.}5\mathinner{.\,.\,.\,.} n}{(n+2)(n+4)\mathinner{.\,.\,.\,.}(2n+1)}F(-\tfrac12(n+1),\ \tfrac12n+1,\ \tfrac12,\ uu)
\]
for odd $n$, with the upper or lower sign according as $\frac12(n+1)$ is even or odd.

The function $U'$ does not admit an equally simple general expression. It is nevertheless easily inferred from the very way the quantities $V,V',V''$, etc., are generated that the last term of $U'$, for even $n$, is
\[
=\pm\frac{2\mathbin{.}2\mathbin{.}4\mathbin{.}4\mathbin{.}6\mathbin{.}6\mathinner{.\,.\,.\,.\,.} n\mathbin{.}n}{3\mathbin{.}5\mathbin{.}7\mathbin{.}9\mathbin{.}11\mathbin{.}13\mathinner{.\,.\,.\,.\,.}(2n-1)(2n+1)}
\]
with the upper or lower sign according as $\frac12n$ is even or odd.

The function $U''=\varphi W^{(n+1)}-V^{(n+1)}$, whose first term we have already shown how to determine in the preceding article, can also be expanded by a recurrence algorithm, since in general we evidently have
\[
\begin{aligned}
\varphi W''-V''&=w'(\varphi W'-V')+v'(\varphi W-V)\\
\varphi W'''-V'''&=w''(\varphi W''-V'')+v''(\varphi W'-V')\\
\varphi W''''-V''''&=w'''(\varphi W'''-V''')+v'''(\varphi W''-V'')
\end{aligned}
\]
\SourceLeaf{198}
etc., and hence, in the case we are considering,
\[
\begin{aligned}&\varphi W^{(m+2)}-V^{(m+2)}=u(\varphi W^{(m+1)}-V^{(m+1)})\\[5pt]&-\frac{(m+1)^2}{(2m-1)(2m+1)}(\varphi W^{(m)}-V^{(m)})\end{aligned}
\]
Thus we find
\[
\begin{aligned}
\varphi W-V&=u^{-1}+\tfrac13u^{-3}+\tfrac15u^{-5}+\tfrac17u^{-7}+\text{etc.}\\
\varphi W'-V'&=\tfrac13u^{-2}+\tfrac15u^{-4}+\tfrac17u^{-6}+\tfrac19u^{-8}+\text{etc.}\\
\varphi W''-V''&=\tfrac4{45}u^{-3}+\tfrac8{105}u^{-5}+\tfrac4{63}u^{-7}+\tfrac{112}{2079}u^{-9}+\text{etc.}\\
\varphi W'''-V'''&=\tfrac4{175}u^{-4}+\tfrac8{315}u^{-6}+\tfrac4{165}u^{-8}+\tfrac{16}{715}u^{-10}+\text{etc.}
\end{aligned}
\]
etc. These series can also be presented as follows:
\[
\begin{aligned}
\varphi W-V&=u^{-1}\left(1+\frac{1\mathbin{.}2}{2\mathbin{.}3}u^{-2}\right.\\[3pt]&\qquad\left.+\frac{1\mathbin{.}2\mathbin{.}3\mathbin{.}4}{2\mathbin{.}4\mathbin{.}3\mathbin{.}5}u^{-4}+\frac{1\mathbin{.}2\mathbin{.}3\mathbin{.}4\mathbin{.}5\mathbin{.}6}{2\mathbin{.}4\mathbin{.}6\mathbin{.}3\mathbin{.}5\mathbin{.}7}u^{-6}+\text{etc.}\right)\\[4pt]
\varphi W'-V'&=\tfrac13u^{-2}\left(1+\frac{2\mathbin{.}3}{2\mathbin{.}5}u^{-4}\right.\\[3pt]&\qquad\left.+\frac{2\mathbin{.}3\mathbin{.}4\mathbin{.}5}{2\mathbin{.}4\mathbin{.}5\mathbin{.}7}u^{-4}+\frac{2\mathbin{.}3\mathbin{.}4\mathbin{.}5\mathbin{.}6\mathbin{.}7}{2\mathbin{.}4\mathbin{.}6\mathbin{.}5\mathbin{.}7\mathbin{.}9}u^{-6}+\text{etc.}\right)\\[4pt]
\varphi W''-V''&=\tfrac4{45}u^{-3}\left(1+\frac{3\mathbin{.}4}{2\mathbin{.}7}u^{-2}\right.\\[3pt]&\qquad\left.+\frac{3\mathbin{.}4\mathbin{.}5\mathbin{.}6}{2\mathbin{.}4\mathbin{.}7\mathbin{.}9}u^{-4}+\frac{3\mathbin{.}4\mathbin{.}5\mathbin{.}6\mathbin{.}7\mathbin{.}8}{2\mathbin{.}4\mathbin{.}6\mathbin{.}7\mathbin{.}9\mathbin{.}11}u^{-6}+\text{etc.}\right)\\[4pt]
\varphi W'''-V'''&=\tfrac4{175}u^{-4}\left(1+\frac{4\mathbin{.}5}{2\mathbin{.}9}u^{-2}\right.\\[3pt]&\qquad\left.+\frac{4\mathbin{.}5\mathbin{.}6\mathbin{.}7}{2\mathbin{.}4\mathbin{.}9\mathbin{.}11}u^{-4}+\frac{4\mathbin{.}5\mathbin{.}6\mathbin{.}7\mathbin{.}8\mathbin{.}9}{2\mathbin{.}4\mathbin{.}6\mathbin{.}9\mathbin{.}11\mathbin{.}13}u^{-6}+\text{etc.}\right)
\end{aligned}
\]
\Needspace{19\baselineskip}
etc. Following this induction, we shall have in general
\[
U''=\varphi W^{(n+1)}-V^{(n+1)}\quad\text{equal to the product of}
\]
\[
\frac{2\mathbin{.}2\mathbin{.}3\mathbin{.}3\mathbin{.}4\mathbin{.}4\mathinner{.\,.\,.\,.\,.}(n+1)\mathbin{.}(n+1)}{3\mathbin{.}3\mathbin{.}5\mathbin{.}5\mathbin{.}7\mathbin{.}7\mathbin{.}9\mathinner{.\,.\,.\,.\,.}(2n+1)(2n+3)}u^{-(n+2)}
\]
and the infinite series
\[
1+\frac{(n+2)(n+3)}{2(2n+5)}u^{-2}+\frac{(n+2)(n+3)(n+4)(n+5)}{2\mathbin{.}4\mathbin{.}(2n+5)(2n+7)}u^{-4}+\text{etc.}
\]
or, if one prefers, $F(\frac12n+1,\ \frac12n+\frac32,\ n+\frac52,\ u^{-2})$. This induction too can very easily be established with complete certainty, either by the well-known method or with the aid of formula 19 in the paper cited several times.

